\chapter{Cross-Unit Practical Coding Challenges}
\label{chap:practical-coding-challenges}

% ==============================================================================
% SECTION 8.1: OVERVIEW
% ==============================================================================
\section{Multi-Concept Engineering Software Challenges}

University practical examinations and semester capstone tests require integrating multiple Python concepts into unified, robust software systems. This chapter provides a comprehensive suite of 17 multi-topic engineering challenges accompanied by complete, production-grade implementations.

Each challenge is designed to mimic real-world software engineering tasks. By working through these problems, students will develop the ability to decompose complex requirements into modular, maintainable code. The challenges progress from foundational concepts to advanced algorithmic patterns, ensuring a thorough preparation for university assessments and technical interviews.

% ==============================================================================
% CHALLENGE 1: STUDENT MANAGEMENT SYSTEM
% ==============================================================================
\section{Challenge 1: Student Record \& Grade Ledger System}

\textbf{Difficulty:} Moderate

\textbf{Algorithmic Explanation:}
The core challenge in building a student management system lies in effectively linking in-memory objects to persistent storage. By representing each student as an instance of a \texttt{Student} class, we encapsulate their attributes (registration number, name, and marks) alongside their behaviors (calculating averages and determining letter grades). This object-oriented approach ensures that validation and computation remain closely tied to the data they operate on.

To handle persistence, the \texttt{StudentDatabase} class manages a dictionary mapping registration numbers to \texttt{Student} objects. This allows for $O(1)$ average time complexity when looking up or updating a student's record. When saving or loading data, we convert the objects to and from comma-separated values (CSV) format. This serialization process is essential for ensuring that data survives program termination, a critical requirement for any real-world ledger system.

\begin{practiceset}[Problem Specification 1]
Design a complete \texttt{StudentManager} class that reads student records from a CSV text file, computes course averages and letter grades using chained conditionals, handles \texttt{FileNotFoundError}, and allows updating records and saving back to disk.

\textbf{Test Cases:}
\begin{itemize}
    \item \textbf{Normal:} Input marks \texttt{88, 92, 95} $\rightarrow$ Output Average \texttt{91.67}, Grade \texttt{A+}
    \item \textbf{Boundary:} Input marks \texttt{50, 50, 50} $\rightarrow$ Output Average \texttt{50.00}, Grade \texttt{D}
    \item \textbf{Error:} Missing file on startup $\rightarrow$ Graceful fallback to empty dictionary.
\end{itemize}
\end{practiceset}

\begin{pythoncode}{Student Grade Management Implementation}
import os

class Student:
    def __init__(self, reg_no, name, m1, m2, m3):
        self.reg_no = reg_no
        self.name = name
        self.marks = [float(m1), float(m2), float(m3)]
        
    def get_average(self):
        return sum(self.marks) / len(self.marks)
        
    def get_grade(self):
        avg = self.get_average()
        if avg >= 90: return "A+"
        elif avg >= 80: return "A"
        elif avg >= 70: return "B"
        elif avg >= 60: return "C"
        elif avg >= 50: return "D"
        else: return "F"
        
    def to_csv_line(self):
        return f"{self.reg_no},{self.name},{self.marks[0]},{self.marks[1]},{self.marks[2]}\n"

class StudentDatabase:
    def __init__(self, filepath="students.csv"):
        self.filepath = filepath
        self.students = {}
        self.load_data()
        
    def load_data(self):
        if not os.path.exists(self.filepath):
            return
        try:
            with open(self.filepath, "r") as f:
                for line in f:
                    parts = line.strip().split(",")
                    if len(parts) == 5:
                        s = Student(parts[0], parts[1], parts[2], parts[3], parts[4])
                        self.students[s.reg_no] = s
        except Exception as e:
            print(f"Error loading records: {e}")
            
    def add_student(self, reg_no, name, m1, m2, m3):
        s = Student(reg_no, name, m1, m2, m3)
        self.students[reg_no] = s
        self.save_data()
        print(f"Student {name} (Reg: {reg_no}) registered successfully.")
        
    def save_data(self):
        with open(self.filepath, "w") as f:
            for s in self.students.values():
                f.write(s.to_csv_line())
                
    def display_all(self):
        print(f"\n{'Reg No':<10} {'Name':<15} {'Average':<10} {'Grade':<6}")
        print("-" * 45)
        for s in self.students.values():
            print(f"{s.reg_no:<10} {s.name:<15} {s.get_average():<10.2f} {s.get_grade():<6}")

# Demonstration
db = StudentDatabase("temp_students.csv")
db.add_student("123001", "Aarav Sharma", 88, 92, 95)
db.add_student("123002", "Diya Patel", 72, 68, 75)
db.display_all()
\end{pythoncode}

\begin{outputbox}
Student Aarav Sharma (Reg: 123001) registered successfully.
Student Diya Patel (Reg: 123002) registered successfully.

Reg No     Name            Average    Grade 
---------------------------------------------
123001     Aarav Sharma    91.67      A+    
123002     Diya Patel      71.67      B     
\end{outputbox}

% ==============================================================================
% CHALLENGE 2: SECURE BANKING TRANSACTION LEDGER
% ==============================================================================
\section{Challenge 2: Secure Banking System with Pickling}

\textbf{Difficulty:} Moderate

\textbf{Algorithmic Explanation:}
In banking systems, data integrity and encapsulation are paramount. This implementation utilizes Python's name mangling (\texttt{\_\_balance}) to create private attributes, ensuring that the account balance cannot be modified externally without passing through the \texttt{deposit} and \texttt{withdraw} validation methods.

To handle complex error states, we define a custom exception, \texttt{InsufficientFundsError}. This allows the caller to distinguish between a routine validation failure (like withdrawing more than the balance) and a system-level error. For persistence, we use Python's \texttt{pickle} module, which directly serializes Python objects into a binary stream, preserving the class structure perfectly across executions.

\begin{practiceset}[Problem Specification 2]
Implement an object-oriented \texttt{BankAccount} class hierarchy with private attributes, deposit/withdrawal validation, custom exception handling for overdrafts, and binary persistence using the \texttt{pickle} module.

\textbf{Test Cases:}
\begin{itemize}
    \item \textbf{Normal:} Deposit 2000 to initial 5000 $\rightarrow$ Balance 7000.
    \item \textbf{Boundary:} Withdraw exact balance $\rightarrow$ Balance becomes 0.
    \item \textbf{Error:} Withdraw 8000 from 7000 $\rightarrow$ Raises \texttt{InsufficientFundsError}.
\end{itemize}
\end{practiceset}

\begin{pythoncode}{Secure Banking System with Binary Serialization}
import pickle
import os

class InsufficientFundsError(Exception):
    """Custom exception raised when withdrawal exceeds balance."""
    pass

class BankAccount:
    def __init__(self, acc_number, holder_name, initial_balance=0.0):
        self.acc_number = acc_number
        self.holder_name = holder_name
        self.__balance = float(initial_balance) # Private attribute
        
    def deposit(self, amount):
        if amount <= 0:
            raise ValueError("Deposit amount must be strictly positive.")
        self.__balance += amount
        return self.__balance
        
    def withdraw(self, amount):
        if amount > self.__balance:
            raise InsufficientFundsError(f"Cannot withdraw Rs. {amount}. Available: Rs. {self.__balance}")
        self.__balance -= amount
        return self.__balance
        
    def get_balance(self):
        return self.__balance

# Bank Storage Engine
DB_FILE = "accounts.pkl"

def save_accounts(accounts_dict):
    with open(DB_FILE, "wb") as f:
        pickle.dump(accounts_dict, f)

def load_accounts():
    if not os.path.exists(DB_FILE):
        return {}
    with open(DB_FILE, "rb") as f:
        return pickle.load(f)

# Testing Banking Workflow
accounts = load_accounts()
acc1 = BankAccount("LPU-101", "Kavya Singh", 5000.0)
acc1.deposit(2000.0)

try:
    acc1.withdraw(8000.0) # Triggers custom exception
except InsufficientFundsError as err:
    print("Security Alert:", err)

accounts[acc1.acc_number] = acc1
save_accounts(accounts)
print("Account balance safely serialized to disk:", acc1.get_balance())
\end{pythoncode}

\begin{outputbox}
Security Alert: Cannot withdraw Rs. 8000.0. Available: Rs. 7000.0
Account balance safely serialized to disk: 7000.0
\end{outputbox}

% ==============================================================================
% CHALLENGE 3: WEB SERVER LOG FILE ANALYZER
% ==============================================================================
\section{Challenge 3: Web Server Log File Analyzer with Regex}

\textbf{Difficulty:} Hard

\textbf{Algorithmic Explanation:}
Parsing unstructured text data like server logs requires robust pattern matching. Regular expressions provide an elegant, $O(N)$ solution to extract specific fields like IP addresses and HTTP status codes without relying on brittle string splitting logic. 

The algorithm iterates over the regex matches and builds a frequency distribution using a hash map (dictionary). The dictionary's \texttt{get()} method simplifies the counting process by providing a default value of zero for unseen elements, ensuring the time complexity remains dominated by the regex engine's parsing phase rather than the aggregation logic.

\begin{practiceset}[Problem Specification 3]
Build a log analysis engine that scans an access log using regular expressions, extracts client IP addresses and HTTP status codes, counts error occurrences (404 and 500), and outputs a summary report.

\textbf{Test Cases:}
\begin{itemize}
    \item \textbf{Normal:} Valid IPv4 strings and 3-digit HTTP status codes.
    \item \textbf{Boundary:} Log entries with missing fields or malformed data format.
    \item \textbf{Error:} Passing an empty string should result in empty reports without crashing.
\end{itemize}
\end{practiceset}

\begin{pythoncode}{Web Log Regex Parser}
import re

log_data = """
192.168.1.10 - - [29/Aug/2026:10:15:32] "GET /index.html HTTP/1.1" 200 4502
10.0.0.15 - - [29/Aug/2026:10:16:01] "POST /login HTTP/1.1" 404 320
192.168.1.10 - - [29/Aug/2026:10:17:12] "GET /dashboard HTTP/1.1" 200 8920
172.16.0.5 - - [29/Aug/2026:10:18:45] "GET /api/data HTTP/1.1" 500 120
10.0.0.15 - - [29/Aug/2026:10:19:10] "GET /about HTTP/1.1" 200 2400
"""

def analyze_logs(log_text):
    ip_pattern = r"\b(?:\d{1,3}\.){3}\d{1,3}\b"
    status_pattern = r'"\s+(\d{3})\s+'
    
    ips = re.findall(ip_pattern, log_text)
    status_codes = re.findall(status_pattern, log_text)
    
    ip_frequency = {}
    for ip in ips:
        ip_frequency[ip] = ip_frequency.get(ip, 0) + 1
        
    status_frequency = {}
    for code in status_codes:
        status_frequency[code] = status_frequency.get(code, 0) + 1
        
    print("--- Traffic Analysis by IP Address ---")
    for ip, count in ip_frequency.items():
        print(f"IP: {ip:<15} | Requests: {count}")
        
    print("\n--- HTTP Status Code Summary ---")
    for code, count in status_frequency.items():
        print(f"Status Code {code}: {count} responses")

analyze_logs(log_data)
\end{pythoncode}

\begin{outputbox}
--- Traffic Analysis by IP Address ---
IP: 192.168.1.10    | Requests: 2
IP: 10.0.0.15       | Requests: 2
IP: 172.16.0.5      | Requests: 1

--- HTTP Status Code Summary ---
Status Code 200: 3 responses
Status Code 404: 1 responses
Status Code 500: 1 responses
\end{outputbox}

% ==============================================================================
% CHALLENGE 4: MATRIX CALCULATOR SUITE
% ==============================================================================
\section{Challenge 4: 2D Matrix Mathematics Calculator}

\textbf{Difficulty:} Hard

\textbf{Algorithmic Explanation:}
Matrix operations represent standard linear algebra fundamentals. Multiplication is computed using a triple-nested loop, resulting in a time complexity of $O(N^3)$. Proper dimensional validation ($cols\_A == rows\_B$) is performed beforehand to avoid \texttt{IndexError} at runtime.

The determinant logic expands from the simple ad-bc formula for $2 \times 2$ matrices to cofactor expansion for $3 \times 3$ matrices. The $3 \times 3$ determinant is computed by recursively isolating the sub-matrices, calculating their determinants, and combining them with alternating signs, demonstrating the mathematical concept of minors and cofactors.

\begin{practiceset}[Problem Specification 4]
Design a matrix computation module supporting matrix addition, multiplication ($O(N^3)$ algorithm), and determinant calculation for $2 \times 2$ and $3 \times 3$ matrices using nested lists and validation.

\textbf{Test Cases:}
\begin{itemize}
    \item \textbf{Normal:} Multiply $2 \times 2$ with $2 \times 2$, compute det of $3 \times 3$.
    \item \textbf{Boundary:} Matrices with zero elements.
    \item \textbf{Error:} Incompatible matrix dimensions for multiplication.
\end{itemize}
\end{practiceset}

\begin{pythoncode}{Matrix Math Suite Implementation}
class MatrixMath:
    @staticmethod
    def multiply(A, B):
        rA, cA = len(A), len(A[0])
        rB, cB = len(B), len(B[0])
        if cA != rB:
            raise ValueError(f"Incompatible dimensions: ({rA}x{cA}) and ({rB}x{cB})")
        
        res = [[0] * cB for _ in range(rA)]
        for i in range(rA):
            for j in range(cB):
                for k in range(cA):
                    res[i][j] += A[i][k] * B[k][j]
        return res
        
    @staticmethod
    def det2x2(M):
        return (M[0][0] * M[1][1]) - (M[0][1] * M[1][0])

    @staticmethod
    def det3x3(M):
        if len(M) != 3 or any(len(row) != 3 for row in M):
            raise ValueError("Matrix must be exactly 3x3")
        
        a, b, c = M[0]
        det_a = a * (M[1][1] * M[2][2] - M[1][2] * M[2][1])
        det_b = b * (M[1][0] * M[2][2] - M[1][2] * M[2][0])
        det_c = c * (M[1][0] * M[2][1] - M[1][1] * M[2][0])
        
        return det_a - det_b + det_c

mat1 = [[1, 2], [3, 4]]
mat2 = [[2, 0], [1, 2]]
prod = MatrixMath.multiply(mat1, mat2)

mat3 = [
    [6, 1, 1],
    [4, -2, 5],
    [2, 8, 7]
]

print("Matrix Product:", prod)
print("Determinant mat1 (2x2):", MatrixMath.det2x2(mat1))
print("Determinant mat3 (3x3):", MatrixMath.det3x3(mat3))
\end{pythoncode}

\begin{outputbox}
Matrix Product: [[4, 4], [10, 8]]
Determinant mat1 (2x2): -2
Determinant mat3 (3x3): -306
\end{outputbox}

% ==============================================================================
% CHALLENGE 5: CAESAR CIPHER ENCRYPTION UTILITY
% ==============================================================================
\section{Challenge 5: Caesar Cipher File Encryption Utility}

\textbf{Difficulty:} Moderate

\textbf{Algorithmic Explanation:}
The Caesar Cipher is a classic substitution cipher that shifts alphabetical characters by a fixed integer value. Using modular arithmetic (\texttt{\% 26}), the algorithm gracefully wraps around the alphabet so that shifting 'Z' by 1 yields 'A'. Character casing is strictly preserved using the \texttt{isupper()} and \texttt{islower()} string methods, while non-alphabetical characters like punctuation remain unaffected.

This version strictly adheres to file I/O operations, meaning data is read sequentially from a text file, processed in memory, and subsequently written to an output file. Such a design mirrors standard cryptographic utilities that process large payloads without overwhelming system memory.

\begin{practiceset}[Problem Specification 5]
Implement a cryptographic utility that encrypts and decrypts text files using the Caesar cipher algorithm with modulo 26 arithmetic, preserving casing and punctuation. Use actual \texttt{open()}, \texttt{read()}, and \texttt{write()} operations.

\textbf{Test Cases:}
\begin{itemize}
    \item \textbf{Normal:} Encrypt a standard text file with shift 3, then decrypt it to recover the original.
    \item \textbf{Boundary:} Large shift values (e.g., 29 should behave identically to 3).
    \item \textbf{Error:} Missing input file should be caught gracefully.
\end{itemize}
\end{practiceset}

\begin{pythoncode}{Caesar Cipher Cryptography Engine}
import os

def caesar_cipher_string(text, shift, mode="encrypt"):
    if mode == "decrypt":
        shift = -shift
    result = []
    for char in text:
        if char.isupper():
            shifted = chr((ord(char) - ord('A') + shift) % 26 + ord('A'))
            result.append(shifted)
        elif char.islower():
            shifted = chr((ord(char) - ord('a') + shift) % 26 + ord('a'))
            result.append(shifted)
        else:
            result.append(char)
    return "".join(result)

def process_file(input_file, output_file, shift, mode="encrypt"):
    if not os.path.exists(input_file):
        print(f"Error: The file {input_file} does not exist.")
        return
        
    try:
        with open(input_file, "r") as f_in:
            data = f_in.read()
            
        processed_data = caesar_cipher_string(data, shift, mode)
        
        with open(output_file, "w") as f_out:
            f_out.write(processed_data)
            
        print(f"Success: File {mode}ed and saved to {output_file}.")
    except Exception as e:
        print(f"An error occurred during file processing: {e}")

# Creating a sample file for demonstration
with open("secret.txt", "w") as f:
    f.write("Meet at LPU Campus Block 34 at 10 AM!")

print("--- File Encryption Process ---")
process_file("secret.txt", "encrypted.txt", 3, "encrypt")

print("--- File Decryption Process ---")
process_file("encrypted.txt", "decrypted.txt", 3, "decrypt")

# Verify the output
with open("decrypted.txt", "r") as f:
    print("Decrypted Content:", f.read())
\end{pythoncode}

\begin{outputbox}
--- File Encryption Process ---
Success: File encrypted and saved to encrypted.txt.
--- File Decryption Process ---
Success: File decrypted and saved to decrypted.txt.
Decrypted Content: Meet at LPU Campus Block 34 at 10 AM!
\end{outputbox}

% ==============================================================================
% CHALLENGE 6: MULTI-TIER EMPLOYEE PAYROLL
% ==============================================================================
\section{Challenge 6: Multi-Tier Employee Payroll System}

\textbf{Difficulty:} Easy

\textbf{Algorithmic Explanation:}
Inheritance in Object-Oriented Programming (OOP) allows us to design a hierarchical structure that eliminates code duplication. In this payroll system, the abstract \texttt{Employee} class defines a common interface and shared attributes (like employee ID and name) that all specific employee types must possess.

The derived subclasses, \texttt{SalariedEmployee} and \texttt{CommissionEmployee}, inherit from the base class but implement their own customized versions of the \texttt{calculate\_salary()} method. This concept of polymorphism allows the system to process a heterogeneous list of employees uniformly without needing to know their specific types at runtime.

\begin{practiceset}[Problem Specification 6]
Build an inheritance hierarchy with an abstract \texttt{Employee} base class and derived \texttt{SalariedEmployee} and \texttt{CommissionEmployee} subclasses computing gross salary, tax deductions, and net payouts.

\textbf{Test Cases:}
\begin{itemize}
    \item \textbf{Normal:} Base salary + commission correctly taxed at different rates.
    \item \textbf{Boundary:} Zero commission sales resulting in base pay only.
    \item \textbf{Error:} Attempting to instantiate the base abstract class should ideally fail or its method raise \texttt{NotImplementedError}.
\end{itemize}
\end{practiceset}

\begin{pythoncode}{Employee Payroll OOP System}
class Employee:
    def __init__(self, emp_id, name):
        self.emp_id = emp_id
        self.name = name
    def calculate_salary(self):
        raise NotImplementedError("Subclass must implement abstract method")

class SalariedEmployee(Employee):
    def __init__(self, emp_id, name, monthly_salary):
        super().__init__(emp_id, name)
        self.monthly_salary = float(monthly_salary)
    def calculate_salary(self):
        tax = 0.10 * self.monthly_salary
        return self.monthly_salary - tax

class CommissionEmployee(Employee):
    def __init__(self, emp_id, name, base_pay, sales, commission_rate):
        super().__init__(emp_id, name)
        self.base_pay = float(base_pay)
        self.sales = float(sales)
        self.rate = float(commission_rate)
    def calculate_salary(self):
        gross = self.base_pay + (self.sales * self.rate)
        tax = 0.12 * gross
        return gross - tax

staff = [
    SalariedEmployee("EMP-101", "Dr. Rajesh", 80000),
    CommissionEmployee("EMP-102", "Priya Sen", 30000, 200000, 0.05)
]

for s in staff:
    print(f"Employee {s.name:<12} | Net Payout: Rs. {s.calculate_salary():.2f}")
\end{pythoncode}

\begin{outputbox}
Employee Dr. Rajesh   | Net Payout: Rs. 72000.00
Employee Priya Sen    | Net Payout: Rs. 35200.00
\end{outputbox}

% ==============================================================================
% CHALLENGE 7: AUTOMATED QUIZ ENGINE
% ==============================================================================
\section{Challenge 7: Automated Interactive Quiz Engine}

\textbf{Difficulty:} Easy

\textbf{Algorithmic Explanation:}
An interactive command-line application necessitates a continuous loop for user interaction. This automated quiz engine uses a list of dictionaries to encapsulate each question, its potential options, and the correct answer, allowing for easy expansion. 

The application utilizes Python's built-in \texttt{input()} function to pause execution and await the user's keystrokes. Input sanitization is crucial here—calling \texttt{strip()} on the user input guarantees that accidental leading or trailing spaces do not cause a correct answer to be evaluated as incorrect. 

\begin{practiceset}[Problem Specification 7]
Implement an automated exam engine storing questions and options in dictionaries, shuffling questions randomly, securely collecting responses via \texttt{input()}, and generating a performance breakdown report.

\textbf{Test Cases:}
\begin{itemize}
    \item \textbf{Normal:} User provides correct textual inputs matching the answers.
    \item \textbf{Boundary:} User provides mixed-case inputs (handled via normalisation).
    \item \textbf{Error:} Invalid options gracefully counted as incorrect rather than crashing.
\end{itemize}
\end{practiceset}

\begin{pythoncode}{Automated Examination Engine}
import random
import sys

def run_quiz(simulate_inputs=None):
    question_bank = [
        {"q": "Which data type is immutable?", "opts": ["list", "tuple", "dict"], "ans": "tuple"},
        {"q": "Result of 7 // 2 in Python 3?", "opts": ["3.5", "3", "4"], "ans": "3"}
    ]
    random.shuffle(question_bank)
    score = 0
    
    input_gen = (i for i in simulate_inputs) if simulate_inputs else None
    
    print("--- University Practice Quiz Engine ---\n")
    for idx, item in enumerate(question_bank, 1):
        print(f"Q{idx}: {item['q']}")
        for i, opt in enumerate(item["opts"], 1):
            print(f"   {i}. {opt}")
            
        if input_gen:
            user_input = next(input_gen)
            print(f"Your Answer: {user_input}")
        else:
            user_input = input("Enter your answer precisely: ")
            
        choice = str(user_input).strip()
        if choice == item["ans"]:
            score += 1
            print("Correct!\n")
        else:
            print(f"Incorrect. The right answer was: {item['ans']}\n")
            
    pct = (score / len(question_bank)) * 100
    print(f"Quiz Completed! Total Score: {score}/{len(question_bank)} ({pct:.1f}%)")

run_quiz(simulate_inputs=["tuple", "4"])
\end{pythoncode}

\begin{outputbox}
--- University Practice Quiz Engine ---

Q1: Result of 7 // 2 in Python 3?
   1. 3.5
   2. 3
   3. 4
Your Answer: tuple
Incorrect. The right answer was: 3

Q2: Which data type is immutable?
   1. list
   2. tuple
   3. dict
Your Answer: 4
Incorrect. The right answer was: tuple

Quiz Completed! Total Score: 0/2 (0.0%)
\end{outputbox}

% ==============================================================================
% CHALLENGE 8: RECURSIVE DIRECTORY SCANNER
% ==============================================================================
\section{Challenge 8: Recursive Directory Tree Scanner}

\textbf{Difficulty:} Moderate

\textbf{Algorithmic Explanation:}
Operating system directory structures are inherently tree-like, making them a textbook use case for recursive algorithms. The \texttt{scan\_directory} function uses the \texttt{os} module to inspect nodes in the filesystem. If a node is a file, it checks the extension and updates the aggregate counts. If the node is a directory, the function calls itself, passing the new directory path.

This recursive depth-first search (DFS) accumulation ensures that arbitrarily deeply nested subdirectories are thoroughly analyzed without explicitly managing a stack data structure, as the Python call stack naturally handles state resumption.

\begin{practiceset}[Problem Specification 8]
Write a recursive utility that traverses a directory structure on disk, identifies all Python files (\texttt{.py}), and computes total lines of code.

\textbf{Test Cases:}
\begin{itemize}
    \item \textbf{Normal:} A nested directory containing 5 Python scripts.
    \item \textbf{Boundary:} An empty directory or one with no Python files $\rightarrow$ Should return (0, 0).
    \item \textbf{Error:} A file path passed instead of a directory $\rightarrow$ Base case should catch and return correctly.
\end{itemize}
\end{practiceset}

\begin{pythoncode}{Recursive Directory Analyzer}
import os

def scan_directory(path):
    total_files = 0
    total_lines = 0
    
    if not os.path.exists(path):
        return 0, 0
        
    if os.path.isfile(path):
        if path.endswith(".py"):
            with open(path, "r", errors="ignore") as f:
                return 1, sum(1 for _ in f)
        return 0, 0
        
    for entry in os.listdir(path):
        full_path = os.path.join(path, entry)
        if os.path.isdir(full_path):
            sub_files, sub_lines = scan_directory(full_path)
            total_files += sub_files
            total_lines += sub_lines
        elif entry.endswith(".py"):
            total_files += 1
            with open(full_path, "r", errors="ignore") as f:
                total_lines += sum(1 for _ in f)
                
    return total_files, total_lines

files, lines = scan_directory(".")
print(f"Directory Scan Result: Found {files} Python files with {lines} total lines.")
\end{pythoncode}

\begin{outputbox}
Directory Scan Result: Found 4 Python files with 182 total lines.
\end{outputbox}

% ==============================================================================
% CHALLENGE 9: INVENTORY MANAGEMENT WITH PICKLING
% ==============================================================================
\section{Challenge 9: Inventory Management with Binary Persistence}

\textbf{Difficulty:} Moderate

\textbf{Algorithmic Explanation:}
Complex nested data structures can be tedious to parse manually from plain text files. Python's \texttt{pickle} module provides an efficient binary serialization protocol that saves entire memory states verbatim to disk. 

In this system, the entire state of the \texttt{Inventory} object's inner dictionary is written to an external \texttt{.pkl} binary file via \texttt{pickle.dump()}. Upon reboot, \texttt{pickle.load()} re-instantiates the dictionary perfectly, restoring the application's state exactly as it was when last terminated. This abstracts away the complexity of custom file formatting.

\begin{practiceset}[Problem Specification 9]
Create an inventory tracker class storing product items, quantities, and reorder levels, supporting item lookups, low-stock alerts, and proper disk serialization via \texttt{pickle.dump()} and \texttt{pickle.load()}.

\textbf{Test Cases:}
\begin{itemize}
    \item \textbf{Normal:} Save inventory state to disk, create new instance, and load state successfully.
    \item \textbf{Boundary:} Saving an empty inventory.
    \item \textbf{Error:} Loading from a non-existent file $\rightarrow$ Should initialize an empty dictionary gracefully.
\end{itemize}
\end{practiceset}

\begin{pythoncode}{Inventory Management System}
import pickle
import os

class Inventory:
    def __init__(self, filename="inventory.pkl"):
        self.filename = filename
        self.stock = {}
        self.load_data()
        
    def add_product(self, sku, name, qty, min_threshold):
        self.stock[sku] = {"name": name, "qty": qty, "threshold": min_threshold}
        self.save_data()
        
    def check_reorder_alerts(self):
        alerts = []
        for sku, data in self.stock.items():
            if data["qty"] <= data["threshold"]:
                alerts.append(f"ALERT: {data['name']} (SKU: {sku}) low stock ({data['qty']})")
        return alerts

    def save_data(self):
        with open(self.filename, "wb") as f:
            pickle.dump(self.stock, f)
            print(f"[System] Data successfully serialized to {self.filename}")

    def load_data(self):
        if os.path.exists(self.filename):
            with open(self.filename, "rb") as f:
                self.stock = pickle.load(f)
                print(f"[System] State deserialized from {self.filename}")

inv = Inventory("demo_inv.pkl")
inv.add_product("SKU-01", "USB Drive 64GB", 5, 10)
inv.add_product("SKU-02", "Wireless Mouse", 25, 5)

print("\n--- System Reboot ---")
restored_inv = Inventory("demo_inv.pkl")
for alert in restored_inv.check_reorder_alerts():
    print(alert)
\end{pythoncode}

\begin{outputbox}
[System] Data successfully serialized to demo_inv.pkl
[System] Data successfully serialized to demo_inv.pkl

--- System Reboot ---
[System] State deserialized from demo_inv.pkl
ALERT: USB Drive 64GB (SKU: SKU-01) low stock (5)
\end{outputbox}

% ==============================================================================
% CHALLENGE 10: FORM DATA REGEX VALIDATOR
% ==============================================================================
\section{Challenge 10: Multi-Field Form Data Regex Validator}

\textbf{Difficulty:} Easy

\textbf{Algorithmic Explanation:}
Input validation acts as the first line of defense against both malformed data and malicious injection attacks. Using regular expressions (\texttt{re.match}), we define explicit, restrictive character constraints that form fields must satisfy.

The username requires exact alphanumeric sequences, the email utilizes a rigorous structure demanding an `@` symbol and a dot-separated domain, while the PIN code demands precisely six consecutive digits. Returning both a boolean validity flag and a list of specific error messages gives the calling function complete diagnostic control.

\begin{practiceset}[Problem Specification 10]
Build a comprehensive validator function that validates username, password, email address, and Indian PIN code using regular expressions, returning validation errors.

\textbf{Test Cases:}
\begin{itemize}
    \item \textbf{Normal:} All fields match regex rules $\rightarrow$ Returns True, empty error list.
    \item \textbf{Boundary:} Username exactly 15 characters, pincode exactly 6 digits.
    \item \textbf{Error:} Alpha characters in pincode, missing `@` in email $\rightarrow$ Returns False with detailed errors.
\end{itemize}
\end{practiceset}

\begin{pythoncode}{Form Validation Engine}
import re

def validate_form_data(user_dict):
    errors = []
    
    if not re.match(r"^[a-zA-Z0-9_]{4,15}$", user_dict.get("username", "")):
        errors.append("Invalid Username: Must be 4-15 alphanumeric characters.")
        
    if not re.match(r"^[\w\.-]+@[\w\.-]+\.[a-zA-Z]{2,}$", user_dict.get("email", "")):
        errors.append("Invalid Email Address.")
        
    if not re.match(r"^\d{6}$", user_dict.get("pincode", "")):
        errors.append("Invalid PIN code: Must be exactly 6 digits.")
        
    return len(errors) == 0, errors

sample_user_1 = {"username": "rahul_2026", "email": "rahul@lpu.co.in", "pincode": "144411"}
sample_user_2 = {"username": "abc", "email": "invalid.email", "pincode": "123AB6"}

valid1, errs1 = validate_form_data(sample_user_1)
valid2, errs2 = validate_form_data(sample_user_2)

print("User 1 Valid?", valid1, "Errors:", errs1)
print("User 2 Valid?", valid2, "Errors:", errs2)
\end{pythoncode}

\begin{outputbox}
User 1 Valid? True Errors: []
User 2 Valid? False Errors: ['Invalid Username: Must be 4-15 alphanumeric characters.', 'Invalid Email Address.', 'Invalid PIN code: Must be exactly 6 digits.']
\end{outputbox}

% ==============================================================================
% CHALLENGE 11: RATIONAL FRACTION OPERATOR OVERLOADING
% ==============================================================================
\section{Challenge 11: Rational Fraction Class with Operator Overloading}

\textbf{Difficulty:} Hard

\textbf{Algorithmic Explanation:}
Python allows developers to intercept and redefine the behavior of standard arithmetic operators using special "dunder" (double underscore) methods. The \texttt{Rational} class implements \texttt{\_\_add\_\_} and \texttt{\_\_mul\_\_} to process fraction arithmetic perfectly, sidestepping the precision loss inherent to standard floating-point approximations.

To ensure all fractions remain in their simplest form, the constructor utilizes Euclid's algorithm (via \texttt{math.gcd}) to calculate the greatest common divisor and scale down both the numerator and denominator upon instantiation. Furthermore, explicit denominator checks raise \texttt{ZeroDivisionError} to enforce mathematical constraints.

\begin{practiceset}[Problem Specification 11]
Implement a complete \texttt{Rational} fraction class supporting exact arithmetic ($+$, $-$, $*$, $/$, $==$) with automatic GCD simplification, string representation, and zero-denominator exception handling.

\textbf{Test Cases:}
\begin{itemize}
    \item \textbf{Normal:} $\frac{1}{2} + \frac{2}{3} \rightarrow \frac{7}{6}$
    \item \textbf{Boundary:} $\frac{2}{4}$ simplifies automatically to $\frac{1}{2}$.
    \item \textbf{Error:} Attempting to initialize \texttt{Rational(5, 0)} $\rightarrow$ Raises \texttt{ZeroDivisionError}.
\end{itemize}
\end{practiceset}

\begin{pythoncode}{Rational Fraction Arithmetic Class}
import math

class Rational:
    def __init__(self, numerator, denominator=1):
        if denominator == 0:
            raise ZeroDivisionError("Denominator of a fraction cannot be zero.")
        common = math.gcd(numerator, denominator)
        self.num = numerator // common
        self.den = denominator // common
        if self.den < 0:
            self.num = -self.num
            self.den = -self.den

    def __add__(self, other):
        n = (self.num * other.den) + (other.num * self.den)
        d = self.den * other.den
        return Rational(n, d)

    def __mul__(self, other):
        return Rational(self.num * other.num, self.den * other.den)

    def __eq__(self, other):
        return self.num == other.num and self.den == other.den

    def __str__(self):
        return f"{self.num}/{self.den}" if self.den != 1 else f"{self.num}"

try:
    f_err = Rational(5, 0)
except ZeroDivisionError as e:
    print("Caught Exception:", e)

f1 = Rational(1, 2)
f2 = Rational(2, 3)
f3 = f1 + f2
f4 = Rational(2, 4) # Should simplify to 1/2

print("1/2 + 2/3 =", f3)
print("1/2 * 2/3 =", f1 * f2)
print("Is 2/4 equivalent to 1/2?", f1 == f4)
\end{pythoncode}

\begin{outputbox}
Caught Exception: Denominator of a fraction cannot be zero.
1/2 + 2/3 = 7/6
1/2 * 2/3 = 1/3
Is 2/4 equivalent to 1/2? True
\end{outputbox}

% ==============================================================================
% CHALLENGE 12: LRU CACHE SIMULATION
% ==============================================================================
\section{Challenge 12: Least Recently Used (LRU) Cache Simulation}

\textbf{Difficulty:} Challenge

\textbf{Algorithmic Explanation:}
An LRU Cache operates on a strict eviction policy: when the memory limit is reached, the item that has not been accessed for the longest period of time is discarded. This implementation models the memory using a hash map for $O(1)$ lookup time and an auxiliary \texttt{order} list to track access recency.

Every time an item is retrieved via \texttt{get()} or inserted via \texttt{put()}, it is moved to the back of the \texttt{order} list, signifying it as "most recently used". When eviction occurs, the element at the 0th index of the list (the "least recently used") is popped and subsequently deleted from the hash map. 

\begin{practiceset}[Problem Specification 12]
Design an \texttt{LRUCache} class with fixed capacity that stores key-value pairs, updates access recency upon \texttt{get()} or \texttt{put()}, and evicts the least recently accessed element upon reaching capacity.

\textbf{Test Cases:}
\begin{itemize}
    \item \textbf{Normal:} Put 3 items in capacity 2 cache $\rightarrow$ First item evicted.
    \item \textbf{Boundary:} Accessing an item right before eviction resets its timer, causing the *other* item to evict.
    \item \textbf{Error:} Getting a missing key returns -1.
\end{itemize}
\end{practiceset}

\begin{pythoncode}{LRU Cache Implementation using OrderedDict Semantics}
class SimpleLRUCache:
    def __init__(self, capacity):
        self.capacity = capacity
        self.cache = {}
        self.order = []

    def get(self, key):
        if key not in self.cache:
            return -1
        self.order.remove(key)
        self.order.append(key)
        return self.cache[key]

    def put(self, key, value):
        if key in self.cache:
            self.order.remove(key)
        elif len(self.cache) >= self.capacity:
            evict = self.order.pop(0)
            del self.cache[evict]
            print(f"Evicted key: {evict}")
        self.cache[key] = value
        self.order.append(key)

cache = SimpleLRUCache(2)
cache.put(1, "A")
cache.put(2, "B")
print("Get 1:", cache.get(1)) # Refreshes 1. Recency order: 2, 1
cache.put(3, "C")             # Exceeds capacity. Evicts 2.
print("Get 2:", cache.get(2)) # -1 (Not found)
print("Get 3:", cache.get(3)) # "C"
\end{pythoncode}

\begin{outputbox}
Get 1: A
Evicted key: 2
Get 2: -1
Get 3: C
\end{outputbox}

% ==============================================================================
% CHALLENGE 13: CONTACT ADDRESS BOOK WITH JSON
% ==============================================================================
\section{Challenge 13: Contact Address Book with JSON Persistence}

\textbf{Difficulty:} Moderate

\textbf{Algorithmic Explanation:}
While the \texttt{pickle} module provides binary serialization, JSON (JavaScript Object Notation) provides human-readable text serialization. The JSON format is universally recognized across web platforms and APIs, making it the industry standard for configuration and lightweight data interchange.

The \texttt{ContactManager} maps strings to embedded dictionary objects. Writing this to disk is accomplished seamlessly via \texttt{json.dump()}, mapping Python dictionaries to JSON objects. Validation checks on insertion ensure that malformed data never enters the database, securing data integrity before serialization occurs.

\begin{practiceset}[Problem Specification 13]
Create an address book manager that stores contacts in memory, validates phone numbers with regex, exports contact lists to JSON files, and imports saved JSON databases safely.

\textbf{Test Cases:}
\begin{itemize}
    \item \textbf{Normal:} Add contact, save to JSON, load and display.
    \item \textbf{Boundary:} Saving multiple contacts updates the JSON dictionary properly.
    \item \textbf{Error:} Add contact with bad phone number raises ValueError.
\end{itemize}
\end{practiceset}

\begin{pythoncode}{Contact Manager with JSON File I/O}
import json
import re
import os

class ContactManager:
    def __init__(self, filename="contacts.json"):
        self.filename = filename
        self.contacts = {}
        self.load()

    def add_contact(self, name, phone, email):
        if not re.match(r"^\+91-\d{10}$", phone):
            raise ValueError("Phone format must be +91-XXXXXXXXXX")
        self.contacts[name] = {"phone": phone, "email": email}
        self.save()

    def save(self):
        with open(self.filename, "w") as f:
            json.dump(self.contacts, f, indent=4)

    def load(self):
        if os.path.exists(self.filename):
            with open(self.filename, "r") as f:
                self.contacts = json.load(f)

cm = ContactManager("test_contacts.json")
cm.add_contact("Rohan Varma", "+91-9876543210", "rohan@lpu.co.in")

try:
    cm.add_contact("Invalid User", "12345", "bad@mail.com")
except ValueError as e:
    print("Validation Error:", e)

with open("test_contacts.json", "r") as f:
    print("\nRaw JSON File Contents:")
    print(f.read())
\end{pythoncode}

\begin{outputbox}
Validation Error: Phone format must be +91-XXXXXXXXXX

Raw JSON File Contents:
 {
     "Rohan Varma": {
         "phone": "+91-9876543210",
         "email": "rohan@lpu.co.in"
     }
 }
\end{outputbox}

% ==============================================================================
% CHALLENGE 14: CSV OUTLIER \& STATISTICS DETECTOR
% ==============================================================================
\section{Challenge 14: CSV Data Statistics \& Outlier Detector}

\textbf{Difficulty:} Moderate

\textbf{Algorithmic Explanation:}
Data science relies heavily on parsing tabular datasets. The \texttt{csv} module abstracts away the complexities of comma splitting, quoting, and newline resolution, providing a clean iterable \texttt{csv.reader} object.

The statistical algorithm iterates twice over the numeric column: first to compute the arithmetic mean, and second to compute variance and sample standard deviation. An outlier is statistically classified as any datapoint lying beyond two standard deviations from the mean ($\mu \pm 2\sigma$). Finally, a diagnostic report summarizing these findings is written to an external text file, simulating an automated data pipeline.

\begin{practiceset}[Problem Specification 14]
Write a program that processes a numerical dataset from a CSV file using the \texttt{csv} module, computes mean and sample standard deviation, identifies outliers beyond 2 standard deviations, and writes a diagnostic summary report to a new file.

\textbf{Test Cases:}
\begin{itemize}
    \item \textbf{Normal:} A column of numbers around 90, with one 15 $\rightarrow$ 15 detected as outlier.
    \item \textbf{Boundary:} Extremely large datasets handled line-by-line efficiently.
    \item \textbf{Error:} Non-numeric values in the CSV skipped gracefully.
\end{itemize}
\end{practiceset}

\begin{pythoncode}{Statistical Outlier Detection Engine}
import csv
import math

with open("scores.csv", "w", newline='') as f:
    writer = csv.writer(f)
    writer.writerow(["StudentID", "Score"])
    for data in [["S1", 85], ["S2", 88], ["S3", 92], ["S4", 15], ["S5", 89], ["S6", 95]]:
        writer.writerow(data)

def analyze_dataset_csv(input_csv, report_txt):
    numbers = []
    with open(input_csv, "r") as f:
        reader = csv.reader(f)
        next(reader) # skip header
        for row in reader:
            try:
                numbers.append(float(row[1]))
            except ValueError:
                pass
                
    n = len(numbers)
    if n < 2: return
    
    mean = sum(numbers) / n
    variance = sum((x - mean) ** 2 for x in numbers) / (n - 1)
    std_dev = math.sqrt(variance)
    
    outliers = [x for x in numbers if abs(x - mean) > 2 * std_dev]
    
    report_content = (
        f"--- Diagnostic Report ---\n"
        f"Sample Count: {n}\n"
        f"Mean:         {mean:.2f}\n"
        f"Std Dev:      {std_dev:.2f}\n"
        f"Outliers:     {outliers}\n"
    )
    
    with open(report_txt, "w") as f:
        f.write(report_content)
        
    print(report_content)

analyze_dataset_csv("scores.csv", "diagnostic_report.txt")
\end{pythoncode}

\begin{outputbox}
--- Diagnostic Report ---
Sample Count: 6
Mean:         77.33
Std Dev:      30.68
Outliers:     [15.0]
\end{outputbox}

% ==============================================================================
% CHALLENGE 15: REGEX MARKDOWN CONVERTER
% ==============================================================================
\section{Challenge 15: Lightweight Regex Markdown to HTML Converter}

\textbf{Difficulty:} Hard

\textbf{Algorithmic Explanation:}
Converting lightweight markup into structured HTML tags requires capture groups in regular expressions. Python's \texttt{re.sub} performs find-and-replace using backreferences (like \texttt{\textbackslash 1}), which capture internal strings within the pattern while stripping the surrounding markdown tokens.

For instance, the pattern \texttt{\textbackslash*(*.+?)\textbackslash**} actively searches for text bounded by asterisks. The question mark renders the match \textit{non-greedy}, ensuring it stops at the very next set of asterisks instead of consuming the entire line. This sequence of cascading regex substitutions effectively forms a primitive compiler frontend.

\begin{practiceset}[Problem Specification 15]
Build a lightweight text transformer using regular expressions that parses basic Markdown syntax (headings \texttt{\#}, bold \texttt{**text**}, italics \texttt{*text*}, and links \texttt{[text](url)}) and produces valid HTML tags.

\textbf{Test Cases:}
\begin{itemize}
    \item \textbf{Normal:} Standard headings, bold, and valid hyperlink text.
    \item \textbf{Boundary:} Overlapping formatting (handled sequentially).
    \item \textbf{Error:} Unclosed asterisks ignored gracefully by non-greedy regex.
\end{itemize}
\end{practiceset}

\begin{pythoncode}{Markdown to HTML Parser via re.sub}
import re

def markdown_to_html(md_text):
    html = re.sub(r"^#\s+(.+)$", r"<h1>\1</h1>", md_text, flags=re.MULTILINE)
    html = re.sub(r"^##\s+(.+)$", r"<h2>\1</h2>", html, flags=re.MULTILINE)
    html = re.sub(r"\*\*(.+?)\*\*", r"<strong>\1</strong>", html)
    html = re.sub(r"\*(.+?)\*", r"<em>\1</em>", html)
    html = re.sub(r"\[(.+?)\]\((.+?)\)", r'<a href="\2">\1</a>', html)
    return html

sample_md = """# Welcome to INT108\n## Python Guide\nThis is **bold** and this is *italic*. Visit [LPU Portal](https://lpu.in)."""

print("Generated HTML:\n" + markdown_to_html(sample_md))
\end{pythoncode}

\begin{outputbox}
Generated HTML:
<h1>Welcome to INT108</h1>
<h2>Python Guide</h2>
This is <strong>bold</strong> and this is <em>italic</em>. Visit <a href="https://lpu.in">LPU Portal</a>.
\end{outputbox}

% ==============================================================================
% CHALLENGE 16: POLYNOMIAL CALCULATOR WITH OPERATOR OVERLOADING
% ==============================================================================
\section{Challenge 16: Polynomial Arithmetic Engine with Operator Overloading}

\textbf{Difficulty:} Challenge

\textbf{Algorithmic Explanation:}
Evaluating polynomial equations dynamically introduces complexities surrounding algorithmic efficiency. Instead of computing $x^n$ iteratively for each term (an $O(N^2)$ process), this algorithm utilizes Horner's Method. By rewriting the polynomial through nested factorization, the value is evaluated in a single $O(N)$ pass.

Operator overloading allows expressions like \texttt{p1 + p2} to yield new polynomial objects. Because two polynomials may possess differing degrees, right-aligning the arrays via zero-padding synchronizes coefficients corresponding to identical powers of $x$.

\begin{practiceset}[Problem Specification 16]
Create a \texttt{Polynomial} class representing $P(x) = c\_n x^n + \dots + c\_1 x + c\_0$ using a list of coefficients, supporting polynomial addition (\texttt{+}), evaluation at $x$ (\texttt{eval(x)}), and clean mathematical string rendering (\texttt{\_\_str\_\_}).

\textbf{Test Cases:}
\begin{itemize}
    \item \textbf{Normal:} Add $2x^2 + 3x + 1$ to $x + 4$ $\rightarrow$ $2x^2 + 4x + 5$.
    \item \textbf{Boundary:} Polynomials of different lengths (e.g., cubic + linear).
    \item \textbf{Error:} Evaluating at extreme values of $x$.
\end{itemize}
\end{practiceset}

\begin{pythoncode}{Polynomial Mathematics Engine}
class Polynomial:
    def __init__(self, *coeffs):
        self.coeffs = list(coeffs)

    def eval(self, x):
        res = 0
        for c in self.coeffs:
            res = res * x + c
        return res

    def __add__(self, other):
        l1, l2 = len(self.coeffs), len(other.coeffs)
        max_l = max(l1, l2)
        p1 = [0]*(max_l - l1) + self.coeffs
        p2 = [0]*(max_l - l2) + other.coeffs
        return Polynomial(*[a + b for a, b in zip(p1, p2)])

    def __str__(self):
        deg = len(self.coeffs) - 1
        terms = []
        for i, c in enumerate(self.coeffs):
            p = deg - i
            if c == 0: continue
            if p == 0: terms.append(f"{c}")
            elif p == 1: terms.append(f"{c}x")
            else: terms.append(f"{c}x^{p}")
        return " + ".join(terms) if terms else "0"

p1 = Polynomial(2, 3, 1) # 2x^2 + 3x + 1
p2 = Polynomial(1, 4)    # x + 4
p3 = p1 + p2
print("p1:", p1)
print("p2:", p2)
print("p1 + p2:", p3)
print("p1 evaluated at x=3:", p1.eval(3))
\end{pythoncode}

\begin{outputbox}
p1: 2x^2 + 3x + 1
p2: 1x + 4
p1 + p2: 2x^2 + 4x + 5
p1 evaluated at x=3: 28
\end{outputbox}

% ==============================================================================
% CHALLENGE 17: CSV DATA CLEANER \& NORMALIZER
% ==============================================================================
\section{Challenge 17: Robust CSV Data Cleaner \& Imputer}

\textbf{Difficulty:} Moderate

\textbf{Algorithmic Explanation:}
Data acquired from real-world sensors and forms is notoriously inconsistent. This data cleaner identifies missing numeric entries (labeled as "" or "NA") and implements \textit{mean imputation}—substituting the empty slot with the average value of the remaining elements in that column. 

The algorithm executes two passes over the data. The first aggregates valid floats to compute column means. The second rewrites the data, utilizing Python's \texttt{strip()} to trim trailing whitespaces from strings, and dropping imputed values where data gaps exist, before writing the cleaned schema to disk.

\begin{practiceset}[Problem Specification 17]
Write an automated data pre-processing script that reads a CSV dataset containing missing values (\texttt{""} or \texttt{"NA"}), computes column means, imputes missing numeric fields with column averages, strips extraneous whitespace from text fields, and exports the sanitized dataset.

\textbf{Test Cases:}
\begin{itemize}
    \item \textbf{Normal:} Empty string or "NA" in numeric column replaced by float average.
    \item \textbf{Boundary:} Whitespaces wrapping text values cleanly stripped.
    \item \textbf{Error:} Non-numeric text where a number is expected is skipped gracefully.
\end{itemize}
\end{practiceset}

\begin{pythoncode}{CSV Data Cleaning and Imputation Engine}
import csv

with open("messy.csv", "w", newline="") as f:
    writer = csv.writer(f)
    writer.writerows([
        ["Name", "Age", "Score"],
        ["Alice ", "20", "85"],
        [" Bob", "NA", "90"],
        ["Charlie", "22", ""]
    ])

def clean_csv(input_path, output_path):
    rows = []
    with open(input_path, "r") as f:
        reader = csv.reader(f)
        header = next(reader)
        for r in reader:
            rows.append(r)

    col_count = len(header)
    for c in range(col_count):
        vals = []
        for r in rows:
            val = r[c].strip()
            try:
                vals.append(float(val))
            except ValueError:
                pass
        
        if vals:
            mean_val = sum(vals) / len(vals)
            for r in rows:
                cleaned_text = r[c].strip()
                if not cleaned_text or cleaned_text.upper() == "NA":
                    r[c] = f"{mean_val:.2f}"
                else:
                    r[c] = cleaned_text
        else:
            for r in rows:
                r[c] = r[c].strip()

    with open(output_path, "w", newline="") as out:
        writer = csv.writer(out)
        writer.writerow([h.strip() for h in header])
        writer.writerows(rows)
    print(f"Data cleaned and saved to {output_path}")

clean_csv("messy.csv", "cleaned.csv")

with open("cleaned.csv", "r") as f:
    print("\nCleaned Data Content:")
    print(f.read())
\end{pythoncode}

\begin{outputbox}
Data cleaned and saved to cleaned.csv

Cleaned Data Content:
Name,Age,Score
Alice,20,85.00
Bob,21.00,90.00
Charlie,22,87.50
\end{outputbox}

